\uuid{nmGi}
\titre{Résolution d'équation différentielle du premier ordre}
\auteur{bodin ; croizat ; sacre}
\organisation{exo7}
\licence{CC BY-NC-SA 4.0}
\datecreate{2022-03-08}
\niveau{L1}
\module{Analyse}
\chapitre{Equation différentielle}
\sousChapitre{Résolution d'équation différentielle du premier ordre}
\difficulte{3}
\theme{}
\exo7qcmid{hilisit.equations-differentielles.034}

\contenu{
\begin{question}
\qclassification{225.01}
Soit $f$ la solution de l'équation différentielle $y'+3y=0$ telle que $f'(0)=-6$. Quelles sont les affirmations vraies ?
\begin{answers}
  \good{La courbe représentative de $f$ passe par $A(0,2)$.}
  \bad{La courbe représentative de $f$ passe par $A(0,-6)$.}
  \bad{$f$ est toujours négative.}
  \good{$f$ est une fonction décroissante sur $\Rr$.}
\end{answers}
\begin{explanations}
Comme $f$ est solution de l'équation différentielle, $f'(0)+3f(0)=0$ donc $f(0)=2$ donc la courbe représentative de $f$ passe par le point de coordonnées $(0,2)$ et ne passe pas par celui de coordonnées $(0,-6)$. De plus, $f(x)=2 e^{-3x}$ et $f'(x)=-6 e^{-3x}$ donc $f$ est toujours positive et $f'$ est toujours négative. Par conséquent $f$ est décroissante sur $\Rr$.

\end{explanations}
\end{question}
}
